\section*{الفصل \ref{ch:b1:stress-strain-rheology} --- الإجهاد والانفعال والريولوجيا}

\begin{solution}{exo:b1:stress-strain-rheology:1}
(أ) $\sigma_v = 2700\times9.81\times10\,000 = \qty{2.65e8}{Pa} = \qty{264.9}{MPa}$.
(ب) $P = 1000\times9.81\times10\,000 = \qty{98.1}{MPa}$. (ج) $\sigma_v' =
264.9 - 98.1 = \qty{166.8}{MPa}$، و$\lambda = 98.1/264.9 = 0.370$، وهو
$1000/2700$.
\end{solution}

\begin{solution}{exo:b1:stress-strain-rheology:2}
(أ) $e = (13-10)/10 = 0.30$ و$S = 1.30$. (ب) $e = (7-10)/10 = -0.30$، أي قِصَر
بنسبة \qty{30}{\percent}، و$S = 0.70$. (ج) الزاوية انغلقت بمقدار
$\psi = \ang{30}$، و$\gamma = \tan\ang{30} = 0.577$.
\end{solution}

\begin{solution}{exo:b1:stress-strain-rheology:3}
(أ) المركز \qty{70}{MPa}، ونصف القطر \qty{30}{MPa}. (ب) أكبر إجهاد قصّ هو نصف القطر، \qty{30}{MPa}، على المستويات التي تصنع \ang{45} مع $\sigma_1$ و
$\sigma_3$، وهي تحمل \omterm{def:b1:stress-strain-rheology:stress}{الإجهاد العمودي} للمركز، \qty{70}{MPa}. (ج)
$2\theta = \ang{60}$: $\sigma_n = 70 + 30\cos\ang{60} = \qty{85}{MPa}$ و
$\tau = 30\sin\ang{60} = \qty{26.0}{MPa}$.
\end{solution}

\begin{solution}{exo:b1:stress-strain-rheology:4}
(أ) $\ang{45} - \ang{30}/2 = \ang{30}$ مع $\sigma_1$، على جانبيه، فالمستويان الممكنان يصنعان \ang{60} مع بعضهما. (ب) مع $c' = 20 + r$، يقرأ شرط \cref{thm:b1:stress-strain-rheology:failure-angle}
$r = C\cos\phi + (20 + r)\sin\phi$، فإن
\[
  r = \frac{C\cos\phi + 20\sin\phi}{1 - \sin\phi} = \frac{8.66 + 10}{0.5}
    = \qty{37.3}{MPa},
\]
و$\sigma_1' = 20 + 2r = \qty{94.6}{MPa}$.
\end{solution}

\begin{solution}{exo:b1:stress-strain-rheology:5}
(أ) $c = 40$ و$R = \sqrt{20^2 + 15^2} = \qty{25}{MPa}$: $\sigma_1 =
\qty{65}{MPa}$ و$\sigma_3 = \qty{15}{MPa}$؛ $\tan2\theta_p = 30/40 = 0.75$،
$2\theta_p = \ang{36.87}$، $\theta_p = \ang{18.43}$. (ب) عند هذه الزاوية
$\cos^2\theta = 0.9$، $\sin^2\theta = 0.1$ و$\sin\theta\cos\theta = 0.3$:
$\sigma_n = 60\times0.9 + 2\times15\times0.3 + 20\times0.1 = 54 + 9 + 2 =
\qty{65}{MPa}$، و$\tau = (20-60)\times0.3 + 15\times(0.9-0.1) = -12 + 12 =
0$. (ج) $R = \qty{25}{MPa}$، على المستويات التي تصنع \ang{45} مع الاتجاهات
الرئيسية.
\end{solution}

\begin{solution}{exo:b1:stress-strain-rheology:6}
(أ) $S_{1,2} = \sqrt{1 + 1} \pm 1$: $S_1 = 2.414$ و$S_2 = 0.414$؛
$\tan2\theta' = 2/2 = 1$، $\theta' = \ang{22.5}$. (ب) $S_1S_2 =
(\sqrt2 + 1)(\sqrt2 - 1) = 1$، وهو $\det F$: لا تغيّر في المساحة. (ج) الخطّ على طول
$(1, 1)$ يصير $F(1, 1) = (1 + 2, 1) = (3, 1)$، عند
$\arctan(1/3) = \ang{18.4}$ من اتجاه القصّ، ومطّه
$\sqrt{10}/\sqrt2 = 2.24$. لقد دار متجاوزًا المحور الطويل للقطع الناقص، الذي
يقع عند \ang{22.5}، نحو مستوى القصّ؛ ومطّه أقلّ قليلًا من $S_1$، كما يجب، لأن
لا خطّ يُمطّ أكثر من $S_1$.
\end{solution}

\begin{solution}{exo:b1:stress-strain-rheology:7}
(أ) $S_1/S_2 = 4$ و$S_1S_2 = 1$ يعطيان $S_1 = 2$ و$S_2 = 0.5$؛ وبحسب
\cref{prop:b1:stress-strain-rheology:simple-shear}، $\gamma = S_1 - S_2 =
1.5$. (ب) $\tan2\theta' = 2/1.5 = 1.333$، $2\theta' = \ang{53.1}$: ينبغي أن تقع
المحاور الطويلة عند \ang{26.6} من الجدران. (ج) $1.5\times200 = \qty{300}{m}$.
\end{solution}

\begin{solution}{exo:b1:stress-strain-rheology:8}
$\sigma_v = 2700\times9.81\times3000 = \qty{79.46}{MPa}$. (أ) $P = \qty{29.43}{MPa}$،
$\sigma_n' = 70 - 29.43 = \qty{40.57}{MPa}$، و
\omterm{prop:b1:stress-strain-rheology:byerlee}{قانون بايرلي} يسمح ب$0.85\times40.57 = \qty{34.5}{MPa}$، وهو أكثر من \qty{30}{MPa} المطبّقة: الصدع
يثبت. (ب) $P = 0.6\times79.46 = \qty{47.68}{MPa}$، $\sigma_n' =
\qty{22.32}{MPa}$، والاحتكاك لا يعطي إلا \qty{19.0}{MPa}: فالصدع
ينزلق. (ج) $30 = 0.85\,(70 - 79.46\,\lambda)$ يعطي $\lambda = (70 -
35.29)/79.46 = 0.437$.
\end{solution}

\begin{solution}{exo:b1:stress-strain-rheology:9}
(أ) يتناسب الإجهاد مع $\dot\varepsilon^{1/4}$: $10^{1/4} = 1.78$. (ب) نسبة
المتانتين هي
\[
  \exp\Bigl[\frac{E}{nR}\Bigl(\frac{1}{673.15} - \frac{1}{773.15}\Bigr)\Bigr]
  = \exp\bigl(6705\times\num{1.921e-4}\bigr) = \exp(1.288) = 3.63 .
\]
(ج) تغيّر المعدّل عشرة أضعاف يغيّر مقاومة الزحف بعامل 1.78، وهو ما تفعله الحرارة بتغيّر قدره
$100\times\ln1.78/\ln3.63 \approx
\qty{45}{\degreeCelsius}$ فقط؛ وعلى منحنى حراري قدره نحو \qty{20}{\degreeCelsius}
لكل كيلومتر، ينقل ذلك الانتقال بنحو \qty{2}{km}. وتغيّر
\omterm{def:b1:heat-budget:geotherm}{المنحنى الحراري} يحرّك متساويات الحرارة، ومعها الانتقال، أكثر بكثير.
\end{solution}

\begin{solution}{exo:b1:stress-strain-rheology:10}
(أ) مع $C = 0$، ينزلق الصدع حين $r = \mu c'/\sqrt{1+\mu^2}$
(\cref{thm:b1:stress-strain-rheology:failure-angle}). هنا $c' = S =
\rho g z(1-\lambda)$، فإن $\sigma_1 - \sigma_3 = 2r = 2\mu S/\sqrt{1+\mu^2}$.
(ب) يزداد $S$ بمقدار $2700\times9.81\times1000\times(1 - 1000/2700) =
\qty{16.68}{MPa}$ لكل كيلومتر، و$2\times0.85/1.3124 = 1.295$: فالميل هو
\qty{21.6}{MPa} لكل كيلومتر، بين \num{13.1} (للصدوع العادية) و
\qty{61.3}{MPa} لكل كيلومتر (للصدوع الدفعية). (ج) $2\mu/(k+\mu) = 0.786$،
$2\mu/k = 1.295$ و$2\mu/(k-\mu) = 3.68$، مع $k = 1.3124$: \omterm{def:b1:stress-strain-rheology:principal}{الإجهاد التفاضلي} عند الانزلاق هو 0.79 و\num{1.30} و\num{3.68} مرة \omterm{def:b1:stress-strain-rheology:lithostatic}{الإجهاد الفعّال} الشاقولي.
\end{solution}

\begin{solution}{exo:b1:stress-strain-rheology:11}
(أ) $t = 60\times\num{e6}\times\num{3.156e7} = \qty{1.89e15}{s}$ و
$2\sqrt{\kappa t} = 2\sqrt{\num{1.89e9}} = \qty{87.0}{km}$. (ب)
$\operatorname{erf}x = 600/1350 = 0.444$ من أجل $x = 0.417$: $z = 0.417\times87.0
= \qty{36.3}{km}$؛ $\operatorname{erf}x = 800/1350 = 0.593$ من أجل $x = 0.586$:
$z = \qty{51.0}{km}$. (ج) ينبغي أن تبلغ الزلازل \qtyrange{36}{51}{km}؛
ويضع المغلّف قاعدة الطبقة الهشّة عند \qty{38.2}{km}، ضمن ذلك
المدى وقريبًا من أعلاه، وعند حرارة \qty{641}{\degreeCelsius}، وهي ضمن
المدى المرصود.
\end{solution}

\begin{solution}{exo:b1:stress-strain-rheology:12}
(أ) $r = 45$ و$c' = \qty{65}{MPa}$؛ المسافة من المركز إلى الخطّ
$\tau = 0.85\sigma_n'$ هي $0.85\times65/1.3124 = \qty{42.1}{MPa}$، أقلّ من
نصف القطر: الدائرة تقطع الخطّ، والمستويات على القوس الذي فوقه
تنزلق. (ب) مع $2\theta$ الزاوية من طرف $\sigma_1$، يحتاج الانزلاق إلى
$45\sin2\theta \ge 0.85\,(65 + 45\cos2\theta)$، أي
$45\sin2\theta - 38.25\cos2\theta \ge 55.25$، أو
$59.06\sin(2\theta - \ang{40.36}) \ge 55.25$، حيث $\ang{40.36} = \phi$.
إذن $\sin(2\theta - \phi) \ge 0.9355$، و$2\theta - \phi$ بين \ang{69.31}
و\ang{110.69}، و$2\theta$ بين \ang{109.67} و\ang{151.06}: العمودان
يقعان بين \ang{54.8} و\ang{75.5} من $\sigma_1$، والكسور المنزلقة
بين \ang{14.5} و\ang{35.2} من $\sigma_1$. ووسط المدى،
\ang{24.8}، هو $\ang{45} - \phi/2$، المستوى الوارد في
\cref{thm:b1:stress-strain-rheology:failure-angle}. (ج) لا ينزلق إلا ذلك المستوى
حين تمسّ الدائرة الخطّ للتوّ: $r = \mu(20 + r)/k$، $r = 20\mu/(k -
\mu) = 17/0.4624 = \qty{36.8}{MPa}$، و$\sigma_1' = 20 + 2r =
\qty{93.5}{MPa}$.
\end{solution}

\begin{solution}{pb:b1:stress-strain-rheology:1}
\textbf{1.} $\sigma_v = 2700\times9.81\times5000 = \qty{132.44}{MPa}$.

\textbf{2.} $P = 1000\times9.81\times5000 = \qty{49.05}{MPa}$، و
$\lambda = 49.05/132.44 = 0.370$.

\textbf{3.} $\sigma_1 = \sigma_v = \qty{132.44}{MPa}$، رأسيًّا، و
$\sigma_3 = \qty{75.00}{MPa}$، أفقيًّا شرق--غرب: \omterm{def:b1:layered-earth:layers}{القشرة} تُمطّ شرقًا--غربًا تحت وزنها نفسه، وهي الحالة التي تقود إلى الصدوع العادية، ومنها هذا الصدع.

\textbf{4.} $\sigma_1 - \sigma_3 = \qty{57.44}{MPa}$؛ المركز
$(132.44 + 75.00)/2 = \qty{103.72}{MPa}$؛ نصف القطر \qty{28.72}{MPa}.

\textbf{5.} $\sigma_1' = 132.44 - 49.05 = \qty{83.39}{MPa}$، $\sigma_3' =
75.00 - 49.05 = \qty{25.95}{MPa}$، والمركز $103.72 - 49.05 = \qty{54.67}{MPa}$.

\textbf{6.} $\sigma_{xx} = \qty{75.00}{MPa}$، $\sigma_{yy} = \qty{132.44}{MPa}$
و$\sigma_{xy} = 0$: المحاور هي الاتجاهات الرئيسية، أحدها رأسيّ، فالمستويات العمودية عليها لا تحمل إجهاد قصّ.

\textbf{7.} على طول انحداره نحو الأسفل يسير الصدع على طول المتجهة الوحدة
$(\cos\ang{60}, -\sin\ang{60})$، شرقًا ونزولًا؛ وعموده الصاعد هو
$\vec n = (\sin\ang{60}, \cos\ang{60}) = (0.866, 0.500)$، عند \ang{30} من
$x$. وزاوية عموده مع $\sigma_1$، على طول $y$، جيبُ تمامها
0.500: $\theta = \ang{60}$.

\textbf{8.} مع \ang{30} من $x$: $\sigma_n = 75.00\times0.75 + 132.44\times0.25
= 56.25 + 33.11 = \qty{89.36}{MPa}$، و$\tau = (132.44 - 75.00)\times
\sin\ang{30}\cos\ang{30} = 57.44\times0.4330 = \qty{24.87}{MPa}$.

\textbf{9.} عند $2\theta = \ang{120}$ من طرف $\sigma_1$:
$\sigma_n = 103.72 + 28.72\cos\ang{120} = 103.72 - 14.36 = \qty{89.36}{MPa}$
و$\tau = 28.72\sin\ang{120} = \qty{24.87}{MPa}$.

\textbf{10.} $\sigma_n' = 89.36 - 49.05 = \qty{40.31}{MPa}$.

\textbf{11.} $24.87/40.31 = 0.617$، أدنى من 0.85: الصدع يثبت.

\textbf{12.} $0.85\times40.31 = \qty{34.26}{MPa}$؛ يقصر إجهاد القصّ عنه بمقدار \qty{9.39}{MPa}.

\textbf{13.} للدائرة مركز عند \qty{54.67}{MPa} ونصف قطر
\qty{28.72}{MPa}، من \num{25.95} إلى \qty{83.39}{MPa}؛ وخطّ بايرلي هو
$\tau = 0.85\,\sigma_n'$ (جميع الإجهادات العمودية هنا دون
\qty{200}{MPa})؛ والصدع هو النقطة $(40.31, 24.87)$، عند \ang{120} من
طرف $\sigma_1$.

\textbf{14.} $0.85\times54.67/\sqrt{1 + 0.85^2} = 46.47/1.3124 =
\qty{35.41}{MPa}$، أكبر من نصف القطر، \qty{28.72}{MPa}: الدائرة كلها
تقع تحت الخطّ، ولا يستطيع أي مستوٍ يمرّ بالصخر أن ينزلق.

\textbf{15.} $\phi = \arctan0.85 = \ang{40.36}$: المستوى الذي يصنع
$\ang{45} - \phi/2 = \ang{24.8}$ مع $\sigma_1$ الشاقولي، وينحدر بزاوية
\ang{65.2}.

\textbf{16.} $S = \sigma_v' = \qty{83.39}{MPa}$ و
$\sigma_1 - \sigma_3 = 2\times0.85\times83.39/(1.3124 + 0.85) =
\qty{65.55}{MPa}$؛ والميل يعطي $13.11\times5 = \qty{65.55}{MPa}$.

\textbf{17.} $\sigma_3 = 132.44 - 65.55 = \qty{66.88}{MPa}$.

\textbf{18.} ينزلق الصدع حين
$r\sin\ang{120} = 0.85\,(c' + r\cos\ang{120})$ مع
$c' = (83.39 + \sigma_3')/2$ و$r = (83.39 - \sigma_3')/2$؛ وحلّ المعادلة يعطي
$\sigma_3' = \qty{17.18}{MPa}$، و$\sigma_3 = \qty{66.23}{MPa}$. وهو أدنى
من المستوى الأفضل: فالصدع لا يقع عند النقطة التي تمسّ فيها الدائرة الخطّ أولًا،
لذلك يجب أن تكبر الدائرة أكثر قبل أن تبلغ نقطة الصدع نفسها.

\textbf{19.} كل إجهاد عمودي فعّال، بما فيه $\sigma_1'$ و$\sigma_3'$، ينخفض بمقدار $\Delta P$، وفرقهما لا يتغيّر؛ والإجهادات القصّية لا تتعلّق بضغط المسام إطلاقًا.

\textbf{20.} ينزلق المستوي الأفضل حين تتساوى المسافة مع نصف القطر:
$0.85\,(54.67 - \Delta P)/1.3124 = 28.72$، فإن
$\Delta P = 54.67 - 28.72\times1.3124/0.85 = 54.67 - 44.34 = \qty{10.33}{MPa}$.

\textbf{21.} عمود الماء في البئر يزن كما كان من قبل، فيجب أن تضيف المضخة كامل \qty{10.33}{MPa} عند رأس البئر، أي نحو 103 أضعاف الضغط الجوي. وهذا يعادل عمود ماء ارتفاعه
$\num{10.33e6}/(1000\times9.81) = \qty{1053}{m}$.

\textbf{22.} مع $\mu = 0.6$ لا يستطيع الصدع أن يتحمّل إلا
$0.6\times40.31 = \qty{24.19}{MPa}$، أقلّ من \qty{24.87}{MPa} التي يحملها:
كان سينزلق من قبل. ولأنه ساكن، فاحتكاكه على الأقل $0.617$، أو أن له بعض \omterm{def:b1:stress-strain-rheology:coulomb}{التماسك}.

\textbf{23.} ينزلق المستوي الأفضل توجّهًا حين يبلغ ضغط المسام
$49.05 + 10.33 = \qty{59.38}{MPa}$، أي $\lambda = 59.38/132.44 = 0.448$،
وهو أدنى بكثير من القيمة 1 للضغط الصخري الساكن.

\textbf{24.} $\Delta P = \sigma_n' - \tau/0.85 = 40.309 - 24.870/0.85 = 40.309 -
29.259 = \qty{11.049}{MPa}$، أكبر من \qty{10.33}{MPa} التي تحرّك صدعًا
من أفضل توجّه. وإلى رقم عشري واحد، فارتفاع ضغط المسام الذي يجعل الصدع ينزلق هو
\textbf{\qty[text-series-to-math=true]{11.0}{MPa}}.
\end{solution}
